\section{Results}

We first report the reaction time (RT) analyses, which assess behavioral sensitivity to morphological and lexical variables, followed by the ERP analyses, which examine how these effects unfold over time in the N250 and N400 components. Within each domain, results are shown separately for words and non-words and for models incorporating orthographic sensitivity and general language proficiency as between-subject factors. These parallel analyses provide complementary perspectives on how reading skill modulates both behavioral and neural measures of morphological processing.

\subsection{Response Time Data with participants grouped by Language Proficiency}

	\subsubsection{Words (Language Proficiency $\times$ Family Size $\times$ Base Frequency) }
	
	We found significant main effects of Base Frequency, $F(1, 93.30) = 10.16$, $p = .002$, $\eta_p^2 = .10$, and Family Size, $F(1, 88.08) = 9.06$, $p = .003$, $\eta_p^2 = .09$, indicating that participants responded more quickly to words derived from high-frequency bases ($M = 601$ ms, $SE = 9.8$) than to those from low-frequency bases ($M = 621$ ms, $SE = 10.3$), and to words from large morphological families ($M = 602$ ms, $SE = 10.0$) than to those from small families ($M = 620$ ms, $SE = 10.1$). These two effects were additive; the Base Frequency $\times$ Family Size interaction was not significant, $F(1, 92.22) = 1.01$, $p = .32$, $\eta_p^2 = .01$. 
	
	Language Proficiency did not yield a main effect, $F(1, 64.88) = 0.00$, $p = .990$, $\eta_p^2 < .001$, nor did it interact with any of the item-level factors (all $p > .54$). Thus, differences in language proficiency did not modulate the pattern of item-level frequency effects in word recognition. 
	
	The model explained approximately 1.7\% of the variance attributable to fixed effects ($R^2_{\text{marginal}} = .017$), consistent with the dominance of subject- and item-level variability in lexical decision response times. Likelihood-ratio tests confirmed that random intercepts for both subjects and items improved model fit, $\chi^2 > 170$, $p < 2 $\times$ 10^{-16}$, while inclusion of random slopes did not ($p > .41$).
	
	\subsubsection{Nonwords}
	
	Response times to nonwords were analyzed as a function of morphological complexity, item-level properties (Family Size or Base Frequency), and Language Proficiency. Separate models were fit for nonwords grouped by family size and by base frequency.
	
	\paragraph{Non-word Reaction Times (Language Proficiency $\times$ Family Size)}
	
	The model revealed a significant main effect of Complexity, $F(1, 62.65) = 88.28$, $p < .001$, $\eta_p^2 = .58$, but no main effects of Family Size, $F(1, 94.52) = 0.91$, $p = .344$, $\eta_p^2 = .01$, or Language Proficiency, $F(1, 63.40) = 0.00$, $p = .957$, $\eta_p^2 < .001$. The three-way interaction among Complexity, Family Size, and Language Proficiency was significant, $F(1, 4387.63) = 4.73$, $p = .030$, $\eta_p^2 = .001$, whereas all two-way interactions were nonsignificant ($p > .49$). The model accounted for 1.6\% of the variance explained by fixed effects ($R^2_{\text{marginal}} = .016$) and 46\% of the total variance ($R^2_{\text{conditional}} = .462$). 
	
	Across conditions, responses to complex nonwords ($M = 734$ ms, $SE = 11.5$) were slower than to simple ones ($M = 700$ ms, $SE = 11.3$). Mean RTs did not differ significantly between large-family ($M = 720$ ms) and small-family ($M = 714$ ms) nonwords, nor between high-proficiency ($M = 716$ ms) and low-proficiency ($M = 717$ ms) participants. Follow-up contrasts showed that for high-proficiency participants, the complexity effect was larger for large-family nonwords (≈45 ms) than for small-family nonwords (≈27 ms), whereas low-proficiency participants showed little or no difference. These results indicate that higher language proficiency magnified the processing cost of morphological complexity for nonwords derived from large morphological families, suggesting greater sensitivity to morphological structure among more proficient readers.
	
	\paragraph{Non-word Reaction Times (Language Proficiency $\times$ Base Frequency)}
	
	The model revealed significant main effects of Complexity, $F(1, 62.89) = 90.74$, $p < .001$, $\eta_p^2 = .59$, and Base Frequency, $F(1, 90.71) = 11.45$, $p = .001$, $\eta_p^2 = .11$. Neither Language Proficiency nor any of its interactions were significant ($p > .30$). A marginal Complexity $\times$ Base Frequency interaction, $F(1, 4491.15) = 3.56$, $p = .059$, $\eta_p^2 < .01$, reflected a slightly smaller complexity cost for nonwords derived from low-frequency bases.
	
	Overall, participants responded more slowly to complex nonwords ($M = 734$ ms, $SE = 11.4$) than to simple ones ($M = 699$ ms, $SE = 11.2$), and to high-frequency-base nonwords ($M = 727$ ms, $SE = 11.3$) than to low-frequency-base nonwords ($M = 707$ ms, $SE = 11.9$). The model explained 2.1\% of the fixed-effect variance ($R^2_{\text{marginal}} = .021$) and 46\% of the total variance ($R^2_{\text{conditional}} = .462$). The direction of effects was consistent across proficiency groups, indicating that differences in language proficiency did not modulate the influence of base frequency or morphological complexity on nonword decision latencies.
	
	\subsubsection{Summary of Response Time Results (Language Proficiency)}
	
	Across both words and nonwords, participants with high and low language proficiency showed comparable overall response times, indicating that general language proficiency did not exert a robust effect on lexical decision performance. For words, response times were independently influenced by Base Frequency and Family Size, with faster responses to words derived from high-frequency bases and from large morphological families. For nonwords, robust effects of Morphological Complexity were again observed, with slower responses to complex than to simple items. Only one interaction involving Language Proficiency reached significance: in the family-size model, high-proficiency participants showed a stronger complexity effect for large-family nonwords, suggesting that proficiency modulates sensitivity to morphological family cues when processing unfamiliar forms. 
	
	Together, these results suggest that while item-level lexical variables (Base Frequency, Family Size, and Complexity) robustly shape response times, individual differences in language proficiency exert subtler influences—detectable primarily in how morphological structure interacts with complexity in the processing of novel forms.
	
	\begin{table}[ht]
		\centering
		\caption{Summary of Response Time Results by Language Proficiency}
		\begin{tabular}{lccccc}
			\toprule
			\textbf{Dataset} & \textbf{Factor} & \textbf{$F$ (df)} & \textbf{$p$} & \textbf{$\eta^2_p$} & \textbf{Direction of Effect} \\
			\midrule
			\textit{Words} & Base Frequency & $F(1, 93.30) = 10.16$ & .002 & .10 & high $<$ low (faster) \\
			& Family Size & $F(1, 88.08) = 9.06$ & .003 & .09 & large $<$ small (faster) \\
			& Language Proficiency & $F(1, 64.88) = 0.00$ & .990 & <.001 & ns \\
			\midrule
			\textit{Non-words (Base Frequency)} & Morphological Complexity & $F(1, 62.89) = 90.74$ & $< .001$ & .59 & complex $>$ simple (slower) \\
			& Base Frequency & $F(1, 90.71) = 11.45$ & .001 & .11 & high $>$ low (slower) \\
			& Language Proficiency & $F(1, 63.36) = 0.00$ & .970 & <.001 & ns \\
			\midrule
			\textit{Non-words (Family Size)} & Morphological Complexity & $F(1, 62.65) = 88.28$ & $< .001$ & .58 & complex $>$ simple (slower) \\
			& Family Size & $F(1, 94.52) = 0.91$ & .344 & .01 & ns \\
			& Language Proficiency & $F(1, 63.40) = 0.00$ & .957 & <.001 & ns \\
			& Complexity $\times$ Family Size $\times$ Language Proficiency & $F(1, 4387.63) = 4.73$ & .030 & .001 & sig. three-way \\
			\bottomrule
		\end{tabular}
	\end{table}


\subsection{Response Time Data with participants grouped by Orthographic Sensitivity}
		
	\subsubsection{Words} 
		
		We found a significant main effect of Base Frequency in that participants responded more quickly to words derived from high-frequency bases, ($M = 602$ ms, SE = 9.6) than to those from low-frequency bases ($M = 622$ ms, SE = 10.1), $F(1, 93.84) = 10.22$, $p = .002$, $\eta_p^2 = .10$—an average facilitation of about 20 ms. A separate main effect of Family Size was also observed,$ F(1, 92.69) = 9.12$, $p = .003$, $\eta_p^2 = .09$, reflecting faster responses to words from large morphological families ($M = 603$ ms, SE = 9.9) than to words from small families ($M = 622$ ms, SE = 9.9). These two effects were additive; the Base Frequency $\times$ Family Size interaction was not significant, $F(1, 93.09) = 1.1$, $p = .30$, $\eta_p^2 = .01$. Orthographic Sensitivity showed only a nonsignificant trend toward faster RTs for higher-sensitivity participants, $F(1, 64.87) = 3.84$, $p = .054,$ $\eta_p^2 = .06$. High-OS participants averaged 595 ms (SE = 12.5), whereas Low-OS participants averaged 630 ms (SE = 13.4). No higher-order interactions approached significance. 
			
		The model explained 36\% of the total variance when both fixed and random effects were included (Conditional $R^2 = .36$), while fixed effects alone accounted for about 3\% (Marginal $R^2 = .03$). These values are comparable to those obtained in the simpler random‐intercept model and are consistent with typical lexical‐decision data, in which between‐subject and item variability dominate.
		
	\subsubsection{Nonwords}
	
		Response times to nonwords were analyzed as a function of morphological complexity, item-level characteristics (family size or base frequency), and individual differences in sensitivity (orthographic or semantic). Separate models were fit for nonwords grouped by either family size or base frequency.
	
		\paragraph{Non-word Reaction Times (Orthographic Sensitivity $\times$ Base Frequency)}
	
		The model revealed significant main effects of Complexity, $F(1,62.46)=93.48$, $p<.001$, $\eta^2_p=.60$, Base Frequency, $F(1,91.60)=11.72$, $p<.001$, $\eta^2_p=.11$, and Orthographic Sensitivity, $F(1,63.49)=5.27$, $p=.025$, $\eta^2_p=.08$. The Complexity $\times$ Base Frequency interaction approached significance, $F(1,4492.31)=3.84$, $p=.050$, $\eta^2_p<.01$, whereas no other interactions were reliable ($p>.52$). The model explained 5\% of the fixed-effect variance ($R^2_{marginal}=.05$) and 46\% of the total variance including random effects ($R^2_{conditional}=.46$).
	
		Overall, responses to complex non-words (736 ms) were slower than to simple ones (701 ms), and high-frequency bases elicited slower responses (729 ms) than low-frequency bases (708 ms). Participants with high orthographic sensitivity responded more quickly (695 ms) than those with low sensitivity (743 ms). The marginal Complexity $\times$ Base Frequency interaction  reflected a slightly larger Complexity effect for high-frequency bases; however, the direction of the effect was consistent across frequency conditions, suggesting a broadly uniform influence of morphological complexity.


		\paragraph{Non-word Reaction Times (Orthographic Sensitivity $$\times$$ Family Size)}
	
		The model revealed significant main effects of Complexity, $F(1, 62.51) = 91.00$, $p < .001$, $\eta^2_p = .59$, and Orthographic Sensitivity, $F(1, 63.54) = 5.33$, $p = .024$, $\eta^2_p = .08$, but no main effect of Family Size, $F(1, 95.24) = 1.03$, $p = .313$, $\eta^2_p = .01$. No interactions were significant ($p > .34$). The model accounted for 4.5\% of the fixed-effect variance ($R^2_{\text{marginal}} = .045$) and 46\% of the total variance including random effects ($R^2_{\text{conditional}} = .461$). Likelihood-ratio tests indicated that inclusion of item-level random intercepts significantly improved model fit, $\chi^2(1) = 144.39$, $p < 2 $\times$ 10^{-16}$, whereas random slopes for Complexity or Family Size did not ($p > .38$).
	
		Across conditions, responses to complex non-words (736 ms) were slower than to simple ones (701 ms), and participants with high orthographic sensitivity responded faster (695 ms) than those with low sensitivity (743 ms). Mean RTs did not differ reliably for large (722 ms) versus small (716 ms) morphological families, and none of the interactions involving family size approached significance.

		
	\subsubsection{Summary of Response Time Results (Orthographic Sensitivity)}
		
	Across both words and non-words, participants with high orthographic sensitivity consistently responded more quickly than those with low sensitivity, indicating that greater sensitivity to orthographic structure facilitated lexical decision performance. For word stimuli, response times were independently influenced by both Base Frequency and Family Size: words derived from high-frequency bases and from large morphological families were recognized more rapidly, and these effects were additive. For non-word stimuli, robust effects of Morphological Complexity were observed in both analyses, with slower responses to complex than to simple items. Non-word response times also showed a modest main effect of Orthographic Sensitivity but only limited effects of lexical variables: Base Frequency produced slower responses for high-frequency bases, and Family Size did not reliably influence responses. 
	
	Taken together, these results suggest that orthographic sensitivity modulates overall response speed, while item-level frequency and morphological structure have dissociable influences depending on lexical status—facilitatory for real words but minimal or inhibitory for novel, morphologically structured non-words.
		
	\begin{table}[ht]
		\centering
		\caption{Summary of Response Time Results by Orthographic Sensitivity}
		\begin{tabular}{lccccc}
			\toprule
			\textbf{Dataset} & \textbf{Factor} & \textbf{$F$ (df)} & \textbf{$p$} & \textbf{$\eta^2_p$} & \textbf{Direction of Effect} \\
			\midrule
			\textit{Words} & Base Frequency & $F(1, 93.84) = 10.22$ & .002 & .10 & high $<$ low (faster) \\
			& Family Size & $F(1, 92.69) = 9.12$ & .003 & .09 & large $<$ small (faster) \\
			& Orthographic Sensitivity & $F(1, 64.87) = 3.84$ & .054 & .06 & high $<$ low (trend) \\
			\midrule
			\textit{Non-words (Base Frequency)} & Morphological Complexity & $F(1, 62.46) = 93.48$ & $< .001$ & .60 & complex $>$ simple (slower) \\
			& Base Frequency & $F(1, 91.60) = 11.72$ & $< .001$ & .11 & high $>$ low (slower) \\
			& Orthographic Sensitivity & $F(1, 63.49) = 5.27$ & .025 & .08 & high $<$ low (faster) \\
			\midrule
			\textit{Non-words (Family Size)} & Morphological Complexity & $F(1, 62.51) = 91.00$ & $< .001$ & .59 & complex $>$ simple (slower) \\
			& Family Size & $F(1, 95.24) = 1.03$ & .313 & .01 & ns \\
			& Orthographic Sensitivity & $F(1, 63.54) = 5.33$ & .024 & .08 & high $<$ low (faster) \\
			\bottomrule
		\end{tabular}
	\end{table}


\subsection{ERP data with participants grouped by Language Proficiency}

\subsubsection{Words}
ERP amplitudes elicited by word targets were analyzed using linear mixed-effects models with Language Proficiency (high, low), Family Size (large, small), and Base Frequency (high, low) as fixed effects. The model included random intercepts for participants and electrodes nested within participants, as well as by-participant random slopes for Family Size and Base Frequency. Because our primary theoretical interest concerned how frequency-based expectations are modulated by morphological structure, interactions are described primarily in terms of base-frequency effects.

\paragraph{N250 Component}

The analysis of N250 amplitudes elicited by word targets revealed no significant main effects of Language Proficiency, Family Size, or Base Frequency (all ps > .22). However, there was a robust Family Size $\times$ Base Frequency interaction, $F(1, 1498) = 32.72$, $p < .001$, $\eta_p^2 = .02$, qualified by a significant Language Proficiency $\times$ Family Size $\times$ Base Frequency interaction, $F(1, 1498) = 16.96$, $p < .001$, $\eta_p^2 = .01$. Collapsing across Language Proficiency, follow-up comparisons indicated that the effect of base frequency differed as a function of family size. For large-family words, high-frequency bases elicited significantly more negative N250 amplitudes than low-frequency bases, $t(67.0) = −2.23$,$ p = .029$, $d = −0.36$. In contrast, base frequency had no reliable effect for small-family words, $t(67.0) = 0.79$, $p = .43$. 

To clarify the three-way interaction, we examined base-frequency effects separately within each combination of Language Proficiency and Family Size. These analyses showed that only high-proficiency participants exhibited a reliable base-frequency effect, and this effect was restricted to large-family words. For these participants, large-family words derived from high-frequency bases elicited more negative N250 amplitudes than those derived from low-frequency bases, $t(67) = 2.47$, $p = .016$. No base-frequency effects were observed for small-family words in the high-Sensitivity group.

Interaction contrasts confirmed that the magnitude of the base-frequency effect differed significantly between large- and small-family words for high-proficiency participants ($\Delta = −1.19 \; \mu V$), t(1498) = −7.08, p < .0001, whereas no such difference was present for low-Sensitivity participants ($\Delta = −0.19 \; \mu V$), $p = .27$. Together, these results indicate that semantic sensitivity modulates how morphological family structure shapes early lexical processing, with high-proficiency readers showing enhanced sensitivity to base-frequency information specifically for large-family words.

\paragraph{N400 Component}

The analysis of N400 amplitudes elicited by word targets revealed no significant main effects of Language Proficiency, Family Size, or Base Frequency (all ps > .14). However, the analysis revealed a robust Family Size $\times$ Base Frequency interaction, $F(1, 1498) = 64.68$, $p < .001$, $\eta_p^2 = .04$, which was further qualified by a significant Language Proficiency $\times$ Family Size $\times$ Base Frequency interaction, $F(1, 1498) = 12.12$, $p < .001$, $\eta_p^2 = .008$.  Collapsing across Language Proficiency, follow-up comparisons indicated that the effect of base frequency differed reliably as a function of family size. For small-family words, high-frequency bases elicited significantly more negative N400 amplitudes than low-frequency bases, $t(72.5) = 3.78$, $p < .001$, $d = 0.50$. In contrast, base frequency did not significantly modulate N400 amplitude for large-family words, $t(72.5) = −1.46$, $p = .15$.  To clarify the three-way interaction, base-frequency effects were examined separately within each combination of Language Proficiency and Family Size. These analyses showed that high-proficiency participants exhibited a robust base-frequency effect for small-family words, such that high-frequency bases elicited more negative N400 amplitudes than low-frequency bases, t(72.5) = 3.99, p < .001, but no reliable base-frequency effect was observed for large-family words. For low-proficiency participants responses to low and high base frequency stimuli did not differ as a result of family size (all $p$s $> .16$). Interaction contrasts confirmed that the difference in the base-frequency effect between large- and small-family words was significantly larger for high-proficiency participants than for low-Sensitivity participants ($\Delta = −0.89 \; \mu V$), $t(1498) = −3.48$, $p = .0005$.


\subsubsection{Non Words}

ERP responses to nonword targets were analyzed using an analogous mixed-effects model with Language Proficiency, Family Size, and Complexity (simple, complex) as fixed effects. As with the word analysis, the model included random intercepts for participants and electrodes nested within participants, as well as by-participant random slopes for Family Size and Complexity.

\paragraph{N250 Component}

The analysis of N250 responses to nonword targets revealed no significant main effects of Semantic Sensitivity, Family Size, or Complexity (all ps > .18). However, the analysis revealed a significant Semantic Sensitivity $	imes$ Family Size $	imes$ Complexity interaction, $F(1, 1498) = 4.71$, $p = .030$, $\eta_p^2 = .003$. Follow-up analyses indicated that this interaction was driven by differences in how complexity effects varied across family size as a function of language proficiency. For high-proficiency participants, N250 amplitudes did not differ reliably between complex and simple nonwords, regardless of family size (all $p$s $> .49$). In contrast, low-proficiency participants showed a trend toward greater N250 negativity for complex relative to simple nonwords when the nonwords belonged to large morphological families, $t(64.3) = 1.82$, $p = .073$, but showed no complexity effect for non-words from small morphological families. Interaction contrasts confirmed that the complexity effect differed significantly across sensitivity groups as a function of morphological family size, $t(1498) = −2.17$, $p = .030$. 


\paragraph{N400 Component}

The analysis of N400 responses to nonword targets revealed a significant main effect of Complexity, $F(1, 58) = 4.15$, $p = .046$, $\eta_p^2 = .07$, such that simple nonwords elicited more negative N400 amplitudes than complex nonwords. In addition, there was a significant Family Size $	imes$ Complexity interaction, $F(1, 1498) = 8.08$, $p = .005$, $\eta_p^2 = .005$, which was further qualified by a robust Language Proficiency $	imes$ Family Size $	imes$ Complexity interaction, $F(1, 1498) = 34.43$, $p < .001$, $\eta_p^2 = .02$.
Collapsing across Language Proficiency, follow-up analyses showed that complexity effects were selectively present for nonwords from large morphological families. Large-family simple nonwords elicited significantly more negative N400 amplitudes than large-family complex nonwords, t(64) = −2.61, p = .011, d = −0.51. No complexity effect was observed for small-family nonwords (all ps > .17), and family size did not reliably modulate N400 amplitude within either complexity level.
Interaction contrasts confirmed that the complexity effect differed significantly across family size conditions, t(1498) = 2.84, p = .005, indicating that morphological family size modulates the impact of structural complexity on nonword processing.
To unpack the three-way interaction, complexity effects were examined separately within each combination of Language Proficiency and Family Size. For low-proficiency participants, large-family complex nonwords elicited significantly more negative N400 amplitudes than large-family simple nonwords, t(64) = 2.79, p = .007. In contrast, no complexity effect was observed for small-family nonwords in this group.
For high-Sensitivity participants, the pattern was reversed. High-proficiency readers showed a marginal tendency toward greater N400 negativity for simple than complex nonwords in small families, t(64) = 1.82, p = .073, whereas complexity had no reliable effect for nonwords from large morphological families.
Interaction contrasts confirmed that the difference in the complexity effect between large- and small-family nonwords differed significantly across proficiency groups, t(1498) = −5.87, p < .0001. Specifically, low-proficiency participants showed a stronger complexity effect for nonwords from large morphological families, whereas high-proficiency participants showed a relatively greater complexity effect for small-family nonwords.




\subsubsection{Summary Across Components}

Across both ERP time windows, effects of language proficiency consistently interacted with morphological variables rather than producing direct amplitude shifts.
At the N250, high-proficiency readers displayed more efficient, selective early parsing, with diminished complexity effects, whereas low-proficiency readers showed enhanced sensitivity to morphological structure at this early stage.

At the N400, proficiency effects emerged more robustly: for words, high-proficiency participants exhibited stronger integration of Family Size and Base Frequency, while for nonwords, they suppressed complexity effects that remained pronounced in low-proficiency readers.
Together, these findings indicate that greater language proficiency modulates both early and late stages of morphological processing—attenuating morpho-orthographic demands (N250) and facilitating semantic integration (N400).